\documentclass[11pt,a4paper]{article}
\usepackage[margin=2.4cm]{geometry}
\usepackage{amsmath,amssymb,amsthm,mathtools}
\usepackage{booktabs}
\usepackage{enumitem}
\usepackage{xcolor}
\usepackage{hyperref}
\usepackage{fontspec}
\usepackage{polyglossia}
\setmainlanguage{russian}
\setotherlanguage{english}
\setmainfont{Times New Roman}
\setmonofont{Consolas}[Scale=0.90]

\hypersetup{
  colorlinks=true,
  linkcolor=blue!70!black,
  citecolor=blue!70!black,
  urlcolor=blue!70!black,
  pdftitle={Об одном замечательном свойстве кусочно-линейного ротатора на кубооктаэдре, тождестве Эйлера и замкнутом вычислении интегралов типа Френеля, Якоби, Бесселя и Дирихле},
  pdfauthor={Каратаев Дмитрий Александрович}
}

\theoremstyle{plain}
\newtheorem{theorem}{Теорема}[section]
\newtheorem{lemma}[theorem]{Лемма}
\newtheorem{proposition}[theorem]{Предложение}
\newtheorem{corollary}[theorem]{Следствие}
\theoremstyle{definition}
\newtheorem{definition}[theorem]{Определение}
\newtheorem{remark}[theorem]{Замечание}

\title{\textbf{Об одном замечательном свойстве кусочно-линейного ротатора на кубооктаэдре, тождестве Эйлера и замкнутом вычислении интегралов типа Френеля, Якоби, Бесселя и Дирихле}}
\author{Д.\,А.~Каратаев}
\date{}

\begin{document}
\maketitle

\begin{abstract}
В работе исследуется дискретное тригонометрическое и интегральное исчисление на трёхмерной оболочке кубооктаэдра $O_2(d)\subset\{-1,0,+1\}^d$ ($d=3$, $|V_{O_2}|=12$ вершин, $|\square|=6$ квадратных и $|\triangle|=8$ треугольных граней) и связанном с ней конечном фазовом кольце $\mathbb{Z}_{N_{\mathrm{ring}}}$ ($N_{\mathrm{ring}} = 4\,Q_\square = 512$, где $Q_\square = 2^{|\triangle|-1} = 128$ --- четверть-период). В продолжение кусочно-линейного анализа Томпсона (\texttt{arXiv:1205.3814}), где оставалась открытой задача построения интегрального исчисления без ветвления по квадрантам, показано, что пара кусочно-линейных проекций $\mathcal{E}_d(\varphi) = (\mathcal{C}_d(\varphi),\mathcal{S}_d(\varphi))$ образует точную дискретную экспоненту Эйлера на $512$-угольнике, порождающую фундаментальную пятёрку $\{0,1,i,\pi,e\}$: сдвиг фазы на четверть периода $+Q_\square(d)=+128$ действует во всех $512$ точках кольца как умножение на мнимую единицу $i$, сдвиг на полупериод $+\Phi_{\mathrm{half}}(d)=+256$ тождественно удовлетворяет соотношению Эйлера $\mathcal{E}_d(\varphi+\Phi_{\mathrm{half}}(d))+\mathcal{E}_d(\varphi)=(0,0)$ ($e^{i\pi}+1=0$), а одношаговый разностный вронскиан $\mathcal{W}_d(\varphi)=\mathcal{C}_d(\varphi)\mathcal{S}_d(\varphi+1)-\mathcal{S}_d(\varphi)\mathcal{C}_d(\varphi+1)=\langle \mathcal{E}_d(\varphi+Q_\square(d)),\Delta\mathcal{E}_d(\varphi)\rangle$ тождественно равен константе $Q_\square(d)=128$ во всех $512$ узлах кольца, включая угловые вершины $L^1$-ромба.

Доказано, что двухтактовая симплектическая рекуррентность второго порядка $\varphi_{t+1}-2\varphi_t+\varphi_{t-1}=\Phi$ объединяет элементарный круговой базис Лиувилля ($\Phi=0$) и неэлементарный квадратично-фазовый базис Френеля--Эйри ($\Phi\ne 0$), а двухкольцевая связка при собственном геометрическом модуле кубооктаэдра $k^2 = 1/\|e\|_2^2 = 1/(d-1) = 1/2$ порождает лемнискатическую триаду эллиптических функций Якоби $(\mathrm{sn}_d,\mathrm{cn}_d,\mathrm{dn}_d)$. Это позволяет вычислить в точных целых числах дискретные интегралы дифракции Френеля ($\mathcal{C}_{\mathrm{Fresnel}}=\mathcal{S}_{\mathrm{Fresnel}}=2048$ за счёт погашения $480$ осциллирующих вычетов антисимметрией полупериода $+256$ и выживания $32$ стационарных точек на биссектрисе $\varphi^2\equiv 64\pmod{512}$), полные и неполные эллиптические интегралы II~рода, интегральный синус Дирихле ($\mathrm{Si}_3(128)=128$), интеграл Парсеваля ($2\,796\,288$) и интеграл Бесселя нулевого порядка ($\mathcal{J}_{0,3}(0)=65\,536$, $\mathcal{J}_{0,3}(1)=32\,768$ и точный корень $\mathcal{J}_{0,3}(4)=0$). Наконец, показано, что континуальная окружность и число $\pi$ возникают как предел спрямления $512$-угольника пифагорова ротатора Вейерштрасса (с собственным конечным дефектом периметра $\pi^2/(6\cdot 512^2)\approx 6{,}3\;\text{ppm}$ на фундаментальном кольце), а проекция $12$ вершин кубооктаэдра $O_2(3)$ на координатную ось порождает равновеликое разбиение Архимеда $(4,4,4)$, которое в сумме с центральным узлом веса $12$ путём прямого сложения 4 + 12 = 16 даёт параболические квадратурные веса $(4,16,4)=4\times(1,4,1)$ и точный трёхмерный континуальный интеграл Римана $\iiint_{[-1,1]^3} F\,\mathrm{d}V$. Все определения и теоремы работы полностью формализованы и машинно проверены в системе доказательства теорем \textbf{Lean~4} без неконструктивных аксиом (\texttt{0 sorry}, \texttt{0 Classical.choice}).

\medskip\noindent
\textbf{Ключевые слова:} кусочно-линейная тригонометрия, тождество Эйлера, разностный вронскиан, подстановка Вейерштрасса, кубооктаэдр, интегралы Френеля, лемнискатические функции Якоби, формальная верификация, Lean~4.

\smallskip\noindent
\textbf{УДК:} 517.3, 517.58, 511.5, 519.6 \qquad \textbf{MSC 2020:} 26A06, 33E05, 39A12, 65D30, 68V20.
\end{abstract}

\section{Введение и постановка задачи}
\label{sec:intro}

Классическая тригонометрия в метрике $L^1$ («манхэттенская», или \emph{taxicab}-тригонометрия на ромбе $|x|+|y|=R$) изучалась в работах Краузе~\cite{Krause1986} и Томпсона~\cite{Thompson2011,Thompson2012}. В частности, в статье~\cite{Thompson2012} (\texttt{arXiv:1205.3814}) были исследованы кусочные производные аналогов синуса, косинуса, тангенса и секанса. Однако при построении замкнутого анализа в~\cite{Thompson2012} возникли три принципиальных препятствия:
\begin{enumerate}[leftmargin=*,noitemsep]
\item кусочно-постоянные производные синуса и косинуса терпят разрыв в четырёх вершинах ромба, вследствие чего привычная дифференциальная связь $(\sin\theta)'=\cos\theta$, $(\cos\theta)'=-\sin\theta$ разрушается;
\item для записи ветвей по квадрантам автору~\cite{Thompson2012} пришлось вводить громоздкие комплексные множители $i^{k-2}$, зависящие от чётности номера квадранта $k$;
\item интегральное исчисление и связь с евклидовой ($L^2$) тригонометрией и классическими неэлементарными интегралами (Френеля, Якоби, Бесселя, Дирихле) остались за рамками исследования.
\end{enumerate}

В настоящей работе показано, что указанные препятствия полностью снимаются при переходе к трёхмерной оболочке кубооктаэдра Архимеда $O_2(d) = \{e=(e_x,e_y,e_z)\in\{-1,0,+1\}^3 \mid e_x^2+e_y^2+e_z^2 = d-1\}$ ($d=3$) и связанному с её граневой структурой целочисленному кольцу вычетов $\mathbb{Z}_{N_{\mathrm{ring}}(d)}$:
\[
  |{\triangle}(d)| = 2^d = 8,\qquad
  |\square(d)| = 2d = 6,\qquad
  |V_{O_2}(d)| = 2d(d-1) = 12,
\]
\[
  d_{\mathrm{in}}(\square) = 4,\qquad
  d_{\mathrm{out}}(E)(d) = 2^d - d = 5,\qquad
  \|e\|_2^2 = d - 1 = 2,
\]
\[
  Q_\square(d) = 2^{|{\triangle}(d)|-1} = 128,\qquad
  \Phi_{\mathrm{half}}(d) = 2\,Q_\square(d) = 256,\qquad
  N_{\mathrm{ring}}(d) = 4\,Q_\square(d) = 512,\qquad
  K_P(d) = \frac{N_{\mathrm{ring}}(d)}{|{\triangle}(d)|} = 64.
\]
Все определения и теоремы настоящей статьи полностью формализованы и машинно проверены в самодостаточном модуле \texttt{RotatorCalculusPaper.lean} для системы \textbf{Lean~4}~\cite{Moura2021}. Всюду далее термин \emph{носитель} (\emph{carrier}) употребляется строго в смысле классической теории моделей А.\,И.~Мальцева (\emph{«Алгебраические системы»}~\cite{Malcev1970}) как основное множество сорта в двусортной алгебраической системе $\langle\mathcal{M},\mathcal{T};\,\pi_{M\to T},\,\iota_{T\to M}\rangle$, где \textbf{дискретный $M$-носитель} $\mathcal{M}=\mathbb{Z}_{512}$ есть конечное кольцо фаз на решётке кубооктаэдра с $L^1$-проекциями $(\mathcal{C}_3,\mathcal{S}_3)$, а \textbf{непрерывный $T$-носитель} $\mathcal{T}$ --- континуальный сорт евклидовой окружности $X^2+Y^2=R^2$, поля $\mathbb{Q}\subset\mathbb{R}$ и интегралов Римана.

\section{Дискретная экспонента Эйлера, тождество \texorpdfstring{$e^{i\pi}+1=0$}{e\^{}(i pi)+1=0}, константный вронскиан и подъём Вейерштрасса на 512-угольник}
\label{sec:wronskian-euler}

Для любого $x\in\mathbb{Z}$ определим симметричную свёртку на главное окно $(-\Phi_{\mathrm{half}}(d),\Phi_{\mathrm{half}}(d)]$:
\begin{equation}
  \label{eq:wrap-ring}
  \mathrm{wrap}_d(x) \;:=\; \Phi_{\mathrm{half}}(d) - \bigl((\Phi_{\mathrm{half}}(d) - x) \bmod N_{\mathrm{ring}}(d)\bigr),
\end{equation}
и кусочно-линейные проекции косинуса и синуса без единого условного ветвления:
\begin{equation}
  \label{eq:proj-cs}
  \mathcal{C}_d(\varphi) \;:=\; Q_\square(d) - |\mathrm{wrap}_d(\varphi)|,\qquad
  \mathcal{S}_d(\varphi) \;:=\; \mathcal{C}_d\bigl(\varphi - Q_\square(d)\bigr).
\end{equation}

\begin{definition}[Дискретная экспонента Эйлера и одношаговый разностный вронскиан]
\label{def:euler-exp-wronskian}
Назовём вектор-пару
\[
  \mathcal{E}_d(\varphi) \;:=\; \bigl(\mathcal{C}_d(\varphi),\,\mathcal{S}_d(\varphi)\bigr) \;\in\; \mathbb{Z}^2
\]
\emph{дискретной экспонентой Эйлера} на кольце $\mathbb{Z}_{N_{\mathrm{ring}}(d)}$, а определитель соседних состояний ротатора
\begin{equation}
  \label{eq:rotator-wronskian}
  \mathcal{W}_d(\varphi) \;:=\; \mathcal{C}_d(\varphi)\,\mathcal{S}_d(\varphi+1) - \mathcal{S}_d(\varphi)\,\mathcal{C}_d(\varphi+1)
  \;=\; \mathcal{C}_d(\varphi)\,\Delta\mathcal{S}_d(\varphi) - \mathcal{S}_d(\varphi)\,\Delta\mathcal{C}_d(\varphi)
\end{equation}
--- \emph{одношаговым разностным вронскианом ротатора}, где $\Delta f(\varphi) = f(\varphi+1)-f(\varphi)$.
\end{definition}

Сопоставим каждой точке $\varphi\in\mathbb{Z}_{N_{\mathrm{ring}}(d)}$ рациональный подъём Брахмагупты--Вейерштрасса на евклидову окружность $X^2+Y^2=R^2$:
\begin{equation}
  \label{eq:weierstrass-lift}
  X_d(\varphi) := \mathcal{C}_d(\varphi)^2 - \mathcal{S}_d(\varphi)^2,\qquad
  Y_d(\varphi) := \|e\|_2^2\,\mathcal{C}_d(\varphi)\,\mathcal{S}_d(\varphi),\qquad
  R_d(\varphi) := \mathcal{C}_d(\varphi)^2 + \mathcal{S}_d(\varphi)^2.
\end{equation}

\begin{theorem}[Дискретное тождество Эйлера, универсальный константный вронскиан, подъём Вейерштрасса и формула Ньютона--Лейбница]
\label{thm:main-1}
При $d=3$ во всех $N_{\mathrm{ring}}(d)=512$ точках фазового кольца $\varphi\in\mathbb{Z}_{N_{\mathrm{ring}}(d)}$ выполнены следующие точные тождества (\texttt{thm1\_constant\_wronskian\_weierstrass\_lift\_and\_newton\_leibniz}):
\begin{enumerate}[label=\textbf{(\roman*)},leftmargin=*]
\item \textbf{Рождение фундаментальной пятёрки $\{0,1,i,\pi,e\}$ и тождество Эйлера $e^{i\pi}+1=0$:}
  \begin{itemize}[noitemsep]
  \item $L^1$-норма ротатора постоянна: $|\mathcal{C}_d(\varphi)| + |\mathcal{S}_d(\varphi)| \equiv Q_\square(d)$, причём $\mathcal{E}_d(0) = (Q_\square(d),0)$ соответствует единице~$1$;
  \item сдвиг фазы на четверть периода $+Q_\square(d)$ действует во всех $512$ точках кольца как точное умножение на мнимую единицу~$i$:
    \[
      \mathcal{E}_d\bigl(\varphi + Q_\square(d)\bigr) \;=\; \bigl(-\mathcal{S}_d(\varphi),\; +\mathcal{C}_d(\varphi)\bigr) \;=\; i\cdot\mathcal{E}_d(\varphi);
    \]
  \item двукратный четверть-оборот $+2Q_\square(d)=+\Phi_{\mathrm{half}}(d)$ ($\pi$) удовлетворяет тождеству Эйлера $e^{i(\varphi+\pi)}+e^{i\varphi}=0$ во всех точках кольца:
    \[
      \mathcal{E}_d\bigl(\varphi + \Phi_{\mathrm{half}}(d)\bigr) + \mathcal{E}_d(\varphi) \;=\; (0,0),
      \qquad
      \mathcal{E}_d\bigl(\Phi_{\mathrm{half}}(d)\bigr) + \mathcal{E}_d(0) \;=\; (0,0);
    \]
  \item скалярное произведение вектора $i\mathcal{E}_d(\varphi)=\mathcal{E}_d(\varphi+Q_\square(d))$ с разностной производной $\Delta\mathcal{E}_d(\varphi)$ совпадает с вронскианом $\mathcal{W}_d(\varphi)$ и \textbf{тождественно равно константе $Q_\square(d)$ во всех $512$ узлах кольца (включая вершины излома)}:
    \begin{equation}
      \label{eq:constant-wronskian-all-512}
      \bigl\langle \mathcal{E}_d\bigl(\varphi+Q_\square(d)\bigr),\,\Delta\mathcal{E}_d(\varphi)\bigr\rangle
      \;=\; \mathcal{W}_d(\varphi) \;\equiv\; Q_\square(d) \quad (=128),\qquad
      \sum_{\varphi=0}^{N_{\mathrm{ring}}(d)-1}\mathcal{W}_d(\varphi) = N_{\mathrm{ring}}(d)\,Q_\square(d).
    \end{equation}
  \end{itemize}
\item \textbf{Евклидов $512$-угольник Вейерштрасса и дифференциал дуги без квадратных корней:}
  во всех $N_{\mathrm{ring}}(d)$ точках выполняется пифагорово тождество $X_d(\varphi)^2 + Y_d(\varphi)^2 = R_d(\varphi)^2$, а одношаговое векторное и скалярное произведения на евклидовой окружности выражаются через скалярный шаг $\mathcal{I}_d(\varphi) := \mathcal{C}_d(\varphi)\mathcal{C}_d(\varphi+1) + \mathcal{S}_d(\varphi)\mathcal{S}_d(\varphi+1)$:
  \begin{equation}
    \label{eq:brahmagupta-step}
    X_d(\varphi)Y_d(\varphi+1) - Y_d(\varphi)X_d(\varphi+1) \;=\; \|e\|_2^2\,Q_\square(d)\,\mathcal{I}_d(\varphi),\qquad
    X_d(\varphi)X_d(\varphi+1) + Y_d(\varphi)Y_d(\varphi+1) \;=\; \mathcal{I}_d(\varphi)^2 - Q_\square(d)^2.
  \end{equation}
\item \textbf{Фундаментальная теорема дискретного анализа (формула Ньютона--Лейбница):}
  кусочно-квадратичные первообразные $\widehat{\Phi}_{\mathcal{C}}(d,\Phi) := \sum_{\varphi=0}^{\Phi-1}\mathcal{C}_d(\varphi)$ и $\widehat{\Phi}_{\mathcal{S}}(d,\Phi) := \sum_{\varphi=0}^{\Phi-1}\mathcal{S}_d(\varphi)$ удовлетворяют $\Delta\widehat{\Phi}_{\mathcal{C}}(d,\varphi) = \mathcal{C}_d(\varphi)$, $\Delta\widehat{\Phi}_{\mathcal{S}}(d,\varphi) = \mathcal{S}_d(\varphi)$, обращаются в точный нуль на полном периоде $N_{\mathrm{ring}}(d)$ и дают точную площадь полуволны $\widehat{\Phi}_{\mathcal{S}}(d,\Phi_{\mathrm{half}}(d)) = Q_\square(d)^2 = 16\,384$.
\end{enumerate}
\end{theorem}

\begin{proof}
Аналитически в первом квадранте $\varphi\in[0,Q_\square(d))$ имеем $\mathcal{C}_d(\varphi)=Q_\square(d)-\varphi$, $\mathcal{S}_d(\varphi)=\varphi$, откуда
\[
  \mathcal{W}_d(\varphi) = (Q_\square(d)-\varphi)(\varphi+1) - \varphi(Q_\square(d)-\varphi-1) = Q_\square(d),
\]
где квадратичные члены $-\varphi^2+\varphi^2$ и линейные члены взаимно уничтожаются. В остальных трёх квадрантах равенство $\mathcal{W}_d(\varphi)\equiv Q_\square(d)$ сохраняется в силу инвариантности определителя относительно поворота на $+Q_\square(d)$: $(-\mathcal{S}_d)(\mathcal{C}_d) - (\mathcal{C}_d)(-\mathcal{S}_d) = \mathcal{C}_d\mathcal{S}_d - \mathcal{S}_d\mathcal{C}_d$. Следовательно, разность дискретного тангенса $\mathcal{T}_d(\varphi)=\mathcal{S}_d(\varphi)/\mathcal{C}_d(\varphi)$ имеет точный вид $\Delta\mathcal{T}_d(\varphi) = Q_\square(d)/(\mathcal{C}_d(\varphi)\mathcal{C}_d(\varphi+1))$, что делает суммирование рациональных тригонометрических дробей телескопическим. Полное машинное доказательство всех пунктов теоремы проверяется редукцией ядра Lean~4 (\texttt{rfl}) в \texttt{thm1\_constant\_wronskian\_weierstrass\_lift\_and\_newton\_leibniz}.
\end{proof}

\section{Собственный неэлементарный базис и замкнутое вычисление интегралов Френеля, Якоби, Дирихле, Парсеваля и Бесселя}
\label{sec:integrals}

В классической дифференциальной алгебре Лиувилля интегралы Френеля $\int\cos(\varphi^2)\,\mathrm{d}\varphi$, эллиптические интегралы Лежандра--Якоби, интегральный синус Дирихле $\mathrm{Si}(\Phi)=\int\frac{\sin\varphi}{\varphi}\,\mathrm{d}\varphi$ и интегралы Бесселя $J_0(k)=\frac{1}{2\pi}\int_0^{2\pi}\cos(k\sin\varphi)\,\mathrm{d}\varphi$ не выражаются через элементарные функции одного кругового ротатора. На кольце $\mathbb{Z}_{N_{\mathrm{ring}}(d)}$ ситуация меняется радикально:
\begin{enumerate}[leftmargin=*,noitemsep]
\item Симплектическая двухтактовая разностная рекуррентность второго порядка (схема Верле)
  \begin{equation}
    \label{eq:verlet-parabola}
    \varphi(t+1) - 2\varphi(t) + \varphi(t-1) = \Phi_{\mathrm{accel}}
    \quad\Longrightarrow\quad
    \varphi(t) = \varphi_0 + \omega_0 t + \Phi_{\mathrm{accel}}\,\frac{t(t-1)}{2}
  \end{equation}
  при нулевом ускорении ($\Phi_{\mathrm{accel}}=0$) порождает линейную фазу (элементарный базис Лиувилля $\mathcal{C}_d(\omega_0 t),\mathcal{S}_d(\omega_0 t)$), а при постоянном ненулевом ускорении ($\Phi_{\mathrm{accel}}\ne 0$) --- квадратичную параболу фазы, суммирование вдоль которой непосредственно задаёт первообразные Френеля--Эйри.
\item На паре связанных колец при естественном геометрическом модуле кубооктаэдра
  \[
    k^2 \;=\; \frac{1}{\|e\|_2^2} \;=\; \frac{1}{d-1} \;=\; \frac{1}{2}
  \]
  (знаменитый \emph{лемнискатический модуль Гаусса} кривой с комплексной мультипликацией в кольце $\mathbb{Z}[i]$) триада Якоби $\mathrm{sn}_d(\varphi):=\mathcal{S}_d(\varphi)$, $\mathrm{cn}_d(\varphi):=\mathcal{C}_d(\varphi)$, $2\,\mathrm{dn}_d(\varphi):=\Phi_{\mathrm{half}}(d)-|\mathcal{S}_d(\varphi)|$ удовлетворяет обоим законам сохранения в точных целых числах.
\end{enumerate}

Определим на кольце $\mathbb{Z}_{N_{\mathrm{ring}}(d)}$ пять классических определённых интегралов:
\begin{align}
  \mathcal{C}_{\mathrm{Fresnel}}(d) &:= \sum_{\varphi=0}^{N_{\mathrm{ring}}(d)-1}\mathcal{C}_d(\varphi^2),\qquad
  \mathcal{S}_{\mathrm{Fresnel}}(d) := \sum_{\varphi=0}^{N_{\mathrm{ring}}(d)-1}\mathcal{S}_d(\varphi^2), \label{eq:def-fresnel}\\
  \mathcal{I}_{\mathrm{ell}}(d;a,b) &:= \sum_{\varphi=0}^{N_{\mathrm{ring}}(d)-1}\bigl(a\,|\mathcal{C}_d(\varphi)| + b\,|\mathcal{S}_d(\varphi)|\bigr), \label{eq:def-elliptic}\\
  \mathrm{Si}_d(\Phi) &:= \sum_{\varphi=1}^{\Phi}\mathrm{tdiv}\bigl(\mathcal{S}_d(\varphi),\,\varphi\bigr), \label{eq:def-dirichlet}\\
  \mathcal{I}_{\mathrm{Parseval}}(d) &:= \sum_{\varphi=0}^{N_{\mathrm{ring}}(d)-1}\mathcal{C}_d(\varphi)^2 = \sum_{\varphi=0}^{N_{\mathrm{ring}}(d)-1}\mathcal{S}_d(\varphi)^2, \label{eq:def-parseval}\\
  \mathcal{J}_{0,d}(k) &:= \sum_{\varphi=0}^{N_{\mathrm{ring}}(d)-1}\mathcal{C}_d\bigl(k\,\mathcal{S}_d(\varphi)\bigr). \label{eq:def-bessel}
\end{align}

\begin{theorem}[Замкнутые значения интегралов Френеля, Якоби, Дирихле, Парсеваля и Бесселя]
\label{thm:main-2}
При $d=3$ на кольце $\mathbb{Z}_{N_{\mathrm{ring}}(d)}$ выполнены следующие точные интегральные соотношения (\texttt{thm2\_closed\_form\_fresnel\_jacobi\_dirichlet\_parseval\_bessel\_integrals}):
\begin{enumerate}[label=\textbf{(\roman*)},leftmargin=*]
\item \textbf{Лемнискатическая триада Якоби и эллиптические интегралы:}
  во всех $N_{\mathrm{ring}}(d)$ точках $|\mathrm{sn}_d(\varphi)|+|\mathrm{cn}_d(\varphi)|=Q_\square(d)$ и $2\,\mathrm{dn}_d(\varphi)+|\mathrm{sn}_d(\varphi)|=\Phi_{\mathrm{half}}(d)$, полный эллиптический интеграл на полуосях $(\|e\|_2^2,d)$ равен $\mathcal{I}_{\mathrm{ell}}(d;\|e\|_2^2,d) = \Phi_{\mathrm{half}}(d)\,Q_\square(d)\,d_{\mathrm{out}}(E)(d) = 163\,840$, а неполный интеграл по первому квадранту равен $Q_\square(d)^2 = 16\,384$.
\item \textbf{Точное схлопывание интегралов дифракции Френеля на биссектрисе стационарной фазы:}
  оба полных интеграла Френеля совпадают и равны
  \begin{equation}
    \label{eq:fresnel-2048}
    \mathcal{C}_{\mathrm{Fresnel}}(d) \;=\; \mathcal{S}_{\mathrm{Fresnel}}(d) \;=\; d_{\mathrm{in}}(\square)\cdot N_{\mathrm{ring}}(d) \;=\; 2048,
  \end{equation}
  причём в точке стационарной фазы $\psi = K_P(d) = 64$ ($\pi/4$) выполнено $\mathcal{C}_d(K_P(d)) = \mathcal{S}_d(K_P(d)) = K_P(d)$.
\item \textbf{Интегральный синус Дирихле и энергетический интеграл Парсеваля:}
  \[
    \mathrm{Si}_d\bigl(Q_\square(d)\bigr) = \mathrm{Si}_d\bigl(\Phi_{\mathrm{half}}(d)\bigr) = Q_\square(d) = 128,
  \]
  \[
    |\square(d)|\cdot\mathcal{I}_{\mathrm{Parseval}}(d) \;=\; N_{\mathrm{ring}}(d)\,\bigl(\|e\|_2^2\,Q_\square(d)^2 + 1\bigr)
    \quad\Longrightarrow\quad
    \mathcal{I}_{\mathrm{Parseval}}(3) = \frac{512\times 32\,769}{6} = 2\,796\,288
  \]
  (множитель $2\cdot 128^2+1 = 32\,769 = 3\times 10\,923$ делится нацело на размерность $d=3$).
\item \textbf{Дискретный интеграл Бесселя $\mathcal{J}_{0,d}(k)$ и его точный корень:}
  \[
    \mathcal{J}_{0,d}(0) = N_{\mathrm{ring}}(d)\,Q_\square(d) = 65\,536,\qquad
    \mathcal{J}_{0,d}(1) = \Phi_{\mathrm{half}}(d)\,Q_\square(d) = 32\,768 = \frac{1}{2}\mathcal{J}_{0,d}(0),\qquad
    \mathcal{J}_{0,d}\bigl(d_{\mathrm{in}}(\square)\bigr) = \mathcal{J}_{0,3}(4) = 0.
  \]
\end{enumerate}
\end{theorem}

\begin{proof}
Покажем концептуальный механизм тождества Френеля~\eqref{eq:fresnel-2048} на кольце $\mathbb{Z}_{512}$. Разобьём $512$ аргументов $\varphi\in\{0,\dots,511\}$ по степени делимости на~$2$:
\begin{itemize}[noitemsep,leftmargin=*]
\item Для $256$ нечётных $\varphi=2k+1$ квадраты $\varphi^2\bmod 512$ равномерно (по $4$ раза) пробегают $64$ вычета вида $8m+1$. Поскольку множество $\{8m+1\pmod{512}\}$ инвариантно относительно сдвига на полупериод $+256 = 8\times 32$, а по теореме~\ref{thm:main-1}(i) $\mathcal{E}_3(\psi+256) = -\mathcal{E}_3(\psi)$, сумма по всем нечётным $\varphi$ тождественно обращается в нуль.
\item Аналогично, для $128$ чисел вида $\varphi=2(2k+1)$ ($\varphi^2\equiv 32m+4$) и $64$ чисел вида $\varphi=4(2k+1)$ ($\varphi^2\equiv 128m+16$) образы инвариантны относительно сдвига $+256$, и их вклад равен нулю.
\item Для $32$ чисел, кратных $16$ ($\varphi=16k$), имеем $\varphi^2\equiv 0\pmod{512}$ при чётном $k$ ($16$ точек, $\mathcal{C}_3(0)=+128$) и $\varphi^2\equiv 256\pmod{512}$ при нечётном $k$ ($16$ точек, $\mathcal{C}_3(256)=-128$), что в сумме также даёт $16\times 128 + 16\times(-128)=0$.
\item Единственная выживающая оболочка состоит из $32$ чисел вида $\varphi = 8(2k+1)$, для каждого из которых $\varphi^2 = 64(2k+1)^2 \equiv 64 = K_P(3) \pmod{512}$ (так как $(2k+1)^2\equiv 1\pmod 8$). Все $32$ точки этой оболочки попадают точно в стационарную точку биссектрисы первого квадранта $\psi = 64$ ($\pi/4 = 45^\circ$), где $\mathcal{C}_3(64)=\mathcal{S}_3(64)=64$, что мгновенно даёт $\mathcal{C}_{\mathrm{Fresnel}}(3)=\mathcal{S}_{\mathrm{Fresnel}}(3) = 32\times 64 = 2048$.
\end{itemize}
В континуальном пределе именно эта выживающая стационарная точка биссектрисы $\pi/4$ даёт проекционный множитель кубооктаэдра $\cos(\pi/4) = 1/\sqrt{\|e\|_2^2} = 1/\sqrt{2}$ в классическом интеграле Френеля $\int_0^\infty\cos(x^2)\,\mathrm{d}x = \frac{1}{\sqrt{2}}\frac{\sqrt{\pi}}{2}$.
\end{proof}

\section{Вывод континуальной окружности, числа \texorpdfstring{$\pi$}{pi} и интеграла Римана из спрямления 512-угольника и поясов Архимеда кубооктаэдра}
\label{sec:continuum-emergence}

В рамках рассматриваемого подхода континуальная окружность и континуальный интеграл Римана не постулируются независимо от кольца $\mathbb{Z}_{N_{\mathrm{ring}}(d)}$, а возникают как теоремы о пределе измельчения (спрямления):
\begin{enumerate}[leftmargin=*,noitemsep]
\item На уровне самого кольца $\mathbb{Z}_{512}$ «круг» есть $512$-угольник пифагорова ротатора~\eqref{eq:weierstrass-lift}, все вершины которого лежат на евклидовой окружности $X_d(\varphi)^2+Y_d(\varphi)^2=R_d(\varphi)^2$ в точных целых числах (в частности, при $\varphi=41$: $(P_x,P_y,P_r)=(14\,703,\,10\,496,\,18\,065)$, $14\,703^2+10\,496^2=18\,065^2$). Переход от $n$-угольника к $2n$-угольнику осуществляется рациональным удвоением Вейерштрасса $t\mapsto 2t/(1-t^2)$ без квадратных корней. При этом на самом $512$-угольнике отношение полупериметра к радиусу отличается от континуального предела $\pi$ на собственную величину
  \begin{equation}
    \label{eq:pi-512-defect}
    \frac{\pi - \pi_{512}}{\pi} \;\approx\; \frac{\pi^2}{|\square(d)|\cdot N_{\mathrm{ring}}(d)^2} \;=\; \frac{\pi^2}{6\cdot 512^2} \;=\; \frac{\pi^2}{1\,572\,864} \;\approx\; 6{,}28\times 10^{-6}\quad (6{,}3\;\text{ppm}),
  \end{equation}
  а при $N\to\infty$ ломаная сходится к окружности по метрике Хаусдорфа (стрела прогиба $\to 0$) и по длине (супремум периметров вписанных ломаных $\to 2\pi r$).
\item Для любого трёхмерного поля $F\colon\mathbb{Z}^3\to\mathbb{Z}$ рассмотрим естественный сумматор шара кубооктаэдра $B_1(d)=\{0\}\cup O_2(d)$, уравнивающий вес центрального узла $(0,0,0)$ с суммой по $|V_{O_2}(d)|=12$ вершинам оболочки кубооктаэдра $O_2(d)$:
  \begin{equation}
    \label{eq:cuboct-3d-integrator}
    \widehat{\mathcal{I}}_{O_2(d)}(F) \;:=\; |V_{O_2}(d)|\cdot F(0,0,0) \;+\; \sum_{e\in O_2(d)} F(e_x,e_y,e_z).
  \end{equation}
  В формуле~\eqref{eq:cuboct-3d-integrator} нет ни одного заранее подставленного квадратурного коэффициента Ньютона--Котеса.
\end{enumerate}

\begin{theorem}[Геометрический вывод континуального интеграла Римана и весов $(1,4,1)/6$ из равновеликих поясов Архимеда кубооктаэдра]
\label{thm:main-3}
При $d=3$ на оболочке кубооктаэдра $O_2(d)$ и кольце $\mathbb{Z}_{N_{\mathrm{ring}}(d)}$ выполнены следующие соотношения (\texttt{thm3\_continuum\_integration\_from\_cuboctahedron\_archimedes\_layers}):
\begin{enumerate}[label=\textbf{(\roman*)},leftmargin=*]
\item \textbf{Равновеликое разбиение Архимеда оболочки кубооктаэдра $O_2(d)$:}
  проекция $12$ вершин $O_2(d)=\mathrm{cubicOrbit3D}(d-1)$ на любую координатную ось $e_x\in\{-1,0,+1\}$ разбивает оболочку на три равновеликих сферических пояса по $d_{\mathrm{in}}(\square)=4$ вершины в каждом:
  \[
    \mathrm{cuboctLayerCount}(d,-1) = \mathrm{cuboctLayerCount}(d,0) = \mathrm{cuboctLayerCount}(d,+1) = d_{\mathrm{in}}(\square) = 4.
  \]
  При сложении экваториального пояса $e_x=0$ ($4$ вершины) с центральным узлом веса $|V_{O_2}(d)|=12$ в операторе~\eqref{eq:cuboct-3d-integrator} послойные веса $(-1:0:+1)$ путём прямого сложения 4 + 12 = 16 собираются в
  \[
    \bigl(d_{\mathrm{in}}(\square),\; |V_{O_2}(d)|+d_{\mathrm{in}}(\square),\; d_{\mathrm{in}}(\square)\bigr)
    \;=\; (4,\,16,\,4) \;=\; d_{\mathrm{in}}(\square)\cdot(1,\,d_{\mathrm{in}}(\square),\,1) \;=\; 4\times(1,4,1)
  \]
  с полной суммой весов $\|e\|_2^2\,|V_{O_2}(d)| = d_{\mathrm{in}}(\square)\,|\square(d)| = 24 = 4\times 6$.
\item \textbf{Точное совпадение с континуальным интегралом Римана $\iiint_{[-1,1]^d} F(x,y,z)\,\mathrm{d}V$:}
  для всех степеней $0\le p\le d$ и трёхмерных квадратичных и кубических форм оператор~\eqref{eq:cuboct-3d-integrator} без остатка воспроизводит континуальный интеграл:
  \[
    \frac{1}{\|e\|_2^2\,|V_{O_2}(d)|}\,\widehat{\mathcal{I}}_{O_2(d)}(x^p) \;=\; \frac{1}{2}\int_{-1}^{+1} x^p\,\mathrm{d}x \quad (0\le p\le 3),
    \qquad
    \frac{1}{\|e\|_2^2\,|V_{O_2}(d)|}\,\widehat{\mathcal{I}}_{O_2(d)}(x^2+y^2+z^2) \;=\; \frac{1}{2^3}\iiint_{[-1,1]^3}(x^2+y^2+z^2)\,\mathrm{d}V = 1.
  \]
\item \textbf{Безостаточное интегрирование убывающих факториальных степеней Похгаммера и нижняя граница радиуса $512$-угольника:}
  для убывающих факториалов $m^{\underline{k}} = m(m-1)\cdots(m-k+1)$ выполнено $d_{\mathrm{in}}(\square)\sum_{m=0}^{|V_{O_2}(d)|} m^{\underline{d}} = (|V_{O_2}(d)|+1)^{\underline{d_{\mathrm{in}}(\square)}}$, а радиус пифагорова ротатора во всех $512$ точках ограничен снизу неравенством $\|e\|_2^2\,R_d(\varphi) \ge Q_\square(d)^2$, при этом знаменатель дефекта периметра $512$-угольника равен $|\square(d)|\,N_{\mathrm{ring}}(d)^2 = d\,\|e\|_2^2\,N_{\mathrm{ring}}(d)^{\|e\|_2^2} = 1\,572\,864$.
\item \textbf{Число $\pi$ как монотонная рациональная функция числа сторон $\pi(n)$, тождество куба вронскиана и устный счёт тригонометрических функций без калькулятора:}
  на $n$-угольнике Вейерштрасса с шагом $s = N_{\mathrm{ring}}(d)/n$ вписанная и описанная рациональные аппроксимации
  \[
    \pi_{\mathrm{in}}(n) = \sum_{j=0}^{n/2-1}\frac{I_s(js)\,W_s(js)}{R_d(js)\,R_d((j+1)s)}\in\mathbb{Q},\qquad
    \pi_{\mathrm{out}}(n) = \sum_{j=0}^{n/2-1}\frac{W_s(js)}{I_s(js)}\in\mathbb{Q}
  \]
  в силу тождества Брахмагупты $R_d(\varphi)R_d(\varphi+s)-I_s(\varphi)^2 \equiv W_s(\varphi)^2$ удовлетворяют точному \emph{тождеству куба вронскиана}:
  \begin{equation}
    \label{eq:wronskian-cube-pi-gap}
    \pi_{\mathrm{out}}(n) - \pi_{\mathrm{in}}(n) \;=\; \sum_{j=0}^{n/2-1}\frac{W_s(js)^3}{I_s(js)\,R_d(js)\,R_d((j+1)s)} \;>\; 0,
  \end{equation}
  откуда при каждом удвоении числа сторон $n\mapsto 2n$ ($s\mapsto s/2$) ширина рациональной вилки сжимается в $2^d/\|e\|_2^2 = d_{\mathrm{in}}(\square) = 4$ раза ($\pi_{\mathrm{in}}(8)=2 < \pi_{\mathrm{in}}(16)=\frac{14}{5} < \pi_{\mathrm{in}}(32)=\frac{6486}{2125} < \pi_{\mathrm{out}}(32)=\frac{2209}{693} < \pi_{\mathrm{out}}(16)=\frac{10}{3} < \pi_{\mathrm{out}}(8)=4$), а взвешенная по размерности $d$ оценка Гюйгенса--Снеллиуса $\pi_d(n) = \bigl(\pi_{\mathrm{in}}(n)+(d-1)\pi_{\mathrm{out}}(n)\bigr)/d$ даёт $\pi_3(16) = \frac{142}{45} = 3{,}1555\dots$ и $\pi_3(512) = 3{,}141592667\dots$ ($8$~верных знаков $\pi$ в чистых рациональных числах). При этом рациональный оператор синуса на параболе Верле $P(p,q)=p(\|e\|_2^2 q - p)$ для доли $(p/q)Q_\square(d)$ прямого угла,
  \[
    \sin_{\mathrm{rat}}(d;\,p,q) \;=\; \frac{d_{\mathrm{in}}(\square)\,p(\|e\|_2^2 q - p)}{d_{\mathrm{out}}(E)(d)\,q^2 - p(\|e\|_2^2 q - p)},
  \]
  вычисляет синус и косинус любого угла без калькулятора в два арифметических действия и на канонических долях прямого угла $p/q\in\{1/d^2,\,1/d,\,1/\|e\|_2^2,\,\|e\|_2^2/d,\,1\}$ ($10^\circ,30^\circ,45^\circ,60^\circ,90^\circ$) возвращает фундаментальные структурные дроби кубооктаэдра:
  \[
    \sin_{\mathrm{rat}}(10^\circ) = \frac{17}{97},\quad
    \sin_{\mathrm{rat}}(30^\circ) = \frac{1}{2},\quad
    \sin_{\mathrm{rat}}(45^\circ) = \frac{12}{17},\quad
    \sin_{\mathrm{rat}}(60^\circ) = \frac{32}{37},\quad
    \sin_{\mathrm{rat}}(90^\circ) = 1.
  \]
\item \textbf{Точная сферическая изотропия второго момента $O(|\mathbf{k}|^2)$, закон анизотропии четвёртого момента $O(|\mathbf{k}|^4)$ и трёхмерный бескорневой ротатор углов Эйлера $(\varphi_1,\varphi_2,\varphi_3)$ на кубооктаэдре:}
  для любого волнового вектора $\mathbf{k}=(k_1,k_2,k_3)$ направленные моменты оболочки кубооктаэдра $O_2(d)$ удовлетворяют точным алгебраическим тождествам:
  \begin{equation}
    \label{eq:cuboct-second-fourth-moments}
    \sum_{e\in O_2(d)} (\mathbf{k}\cdot e)^2 \;=\; |{\triangle}(d)|\,|\mathbf{k}|^2 \;=\; 8\,|\mathbf{k}|^2,\qquad
    \sum_{e\in O_2(d)} (\mathbf{k}\cdot e)^4 \;=\; |V_{O_2}(d)|\,|\mathbf{k}|^4 - d_{\mathrm{in}}(\square)\sum_{i=1}^d k_i^4 \;=\; 12\,|\mathbf{k}|^4 - 4\sum_{i=1}^3 k_i^4,
  \end{equation}
  откуда дисперсионное разложение $\sum_{e\in O_2(d)}(1-\cos(\mathbf{k}\cdot e)) = 4\,|\mathbf{k}|^2 - \bigl(\frac{1}{2}|\mathbf{k}|^4 - \frac{1}{6}\sum_{i=1}^3 k_i^4\bigr) + O(|\mathbf{k}|^6)$ строго изотропно в ведущем порядке $O(|\mathbf{k}|^2)$, а в четвёртом порядке $O(|\mathbf{k}|^4)$ даёт нормированный четвёртый момент $12-4=8=|{\triangle}(d)|$ вдоль координатной оси $(1,0,0)$, $12-2=10=2\,d_{\mathrm{out}}(E)(d)$ вдоль диагонали грани $(1,1,0)$ и отношение анизотропного инварианта $\sum_{i=1}^d k_i^4/|\mathbf{k}|^4$ между осью $(1,0,0)$ и пространственной диагональю $(1,1,1)$, равное в точности $d:1 = 3:1$. При этом число независимых плоскостей вращения и наклона (число углов Эйлера $\dim\mathrm{SO}(d)=d(d-1)/2$) совпадает с размерностью пространства $d$ исключительно при $d=3$, разбивая $|V_{O_2}(d)| = \frac{d(d-1)}{2}\,d_{\mathrm{in}}(\square) = 3\times 4 = 12$ вершин кубооктаэдра на три взаимно ортогональных экваториальных квадрата $(xy,xz,yz)$, а композиция трёх плоских пифагоровых поворотов Вейерштрасса $\mathcal{R}_{\mathrm{Euler}}(\varphi_1,\varphi_2,\varphi_3)$ по углам Эйлера $(\varphi_1,\varphi_2,\varphi_3)\in\mathbb{Z}_{N_{\mathrm{ring}}(d)}^3$ осуществляет трёхмерные повороты и наклоны в чистых целых числах с точным сохранением сферической нормы $\|\mathcal{R}_{\mathrm{Euler}}(\varphi_1,\varphi_2,\varphi_3)\mathbf{v}\|_2^2 = \bigl(R_d(\varphi_1)R_d(\varphi_2)R_d(\varphi_3)\bigr)^2\|\mathbf{v}\|_2^2$ и переводит при наклоне $\varphi_2 = K_P(d)=64$ вершину $(1,1,0)$ плоскости $xy$ через треугольную грань кубооктаэдра точно в вершину $(0,1,1)$ плоскости $yz$.
\end{enumerate}
\end{theorem}

\section{Функционально полный булев базис Поста, представление 3-SAT на треугольных гранях \texorpdfstring{$\mathbb{Z}_{512}$}{Z512} и алгебраическая асимметрия прямой проверки и обращения}
\label{sec:boolean-3sat}

На антиподальных полюсах кольца $\{0,\Phi_{\mathrm{half}}(d)\}=\{0,256\}\subset\mathbb{Z}_{N_{\mathrm{ring}}(d)}$ ($\mathcal{C}_3(0)=+128$ кодирует истину~$1$, $\mathcal{C}_3(256)=-128$ кодирует ложь~$0$) кусочно-линейный ротатор индуцирует два алгебраических канала на треугольных гранях $(\psi_1,\psi_2,\psi_3)\in\{0,256\}^3$:
\begin{enumerate}[leftmargin=*,noitemsep]
\item \textbf{Канал~1 (линейный класс Поста~$L$ --- сложение аргументов до оператора $\mathcal{C}_3$):}
  сдвиг $\mathrm{NOT}(\psi)=(\psi+256)\bmod 512$ меняет знак проекции ($\mathcal{C}_3(\psi+256)=-\mathcal{C}_3(\psi)$), а сложение фаз задаёт трехаргументную чётность $\mathrm{XOR}_3$:
  \begin{equation}
    \label{eq:pre-rotator-xor}
    Q_\square(d)^2\,\mathcal{C}_d(\psi_1+\psi_2+\psi_3) \;\equiv\; \mathcal{C}_d(\psi_1)\,\mathcal{C}_d(\psi_2)\,\mathcal{C}_d(\psi_3),
  \end{equation}
  где подгруппа полюсов $(\{0,256\},+\bmod 512)\cong(\mathbb{F}_2,\oplus)$ ($256+256\equiv 0\pmod{512}$).
\item \textbf{Канал~2 (нелинейный пороговый класс --- сложение значений после оператора $\mathcal{C}_3$ с целочисленным делением на $N_{\mathrm{ring}}(d)=512$):}
  положим $\Delta_{\mathrm{3SAT}}(\psi_1,\psi_2,\psi_3) := Q_\square(d)-\bigl(\mathcal{C}_d(\psi_1)+\mathcal{C}_d(\psi_2)+\mathcal{C}_d(\psi_3)\bigr)$. Поскольку при $d=3$ выполнено $(1+d)Q_\square(d)=4\times 128=512=N_{\mathrm{ring}}(d)$, неполные частные деления на $N_{\mathrm{ring}}(d)$ со знаком к нулю задают дизъюнкцию и конъюнкцию без условных ветвлений:
  \begin{equation}
    \label{eq:post-rotator-or-and}
    q_{\mathrm{OR}}(\psi_1,\psi_2,\psi_3) := \mathrm{tdiv}\bigl(\Delta_{\mathrm{3SAT}}(\psi_1,\psi_2,\psi_3),\,N_{\mathrm{ring}}(d)\bigr)\in\{0,1\},\qquad
    q_{\mathrm{AND}}(\psi_1,\psi_2,\psi_3) := \mathrm{tdiv}\Bigl(Q_\square(d)+\sum_{k=1}^3\mathcal{C}_d(\psi_k),\,N_{\mathrm{ring}}(d)\Bigr)\in\{0,1\}.
  \end{equation}
\end{enumerate}

\begin{theorem}[Функциональная полнота Поста \{\texttt{NOT}, \texttt{XOR}, \texttt{OR}, \texttt{AND}\}, представление 3-SAT, склейка остатка $7\to 1$ и некоммутативность двух каналов на $\mathbb{Z}_{512}$]
\label{thm:main-4}
При $d=3$ на полюсах $\{0,256\}^3\subset\mathbb{Z}_{512}^3$ выполнены следующие точные соотношения (\texttt{thm4\_boolean\_post\_completeness\_3sat\_reduction\_and\_channel\_orthogonality}):
\begin{enumerate}[label=\textbf{(\roman*)},leftmargin=*]
\item \textbf{Функциональная полнота Поста и эквивалентность выполнимости 3-SAT нулю частных Канала~2:}
  во всех $2^3=8$ состояниях $\{0,256\}^3$ выполнено мультипликативное тождество чётности~\eqref{eq:pre-rotator-xor}, частное $q_{\mathrm{AND}}$ равно~$1$ исключительно на $(0,0,0)$ и~$0$ на остальных $7$~октантах, а частное дефекта $q_{\mathrm{OR}}$ равно~$0$ на всех $7$~выполняющих октантах дизъюнкта $(l_1\lor l_2\lor l_3)$ и равно~$1$ ($\Delta_{\mathrm{3SAT}}=512$) исключительно на запрещённом октанте $(256,256,256)$. Тем самым в силу полиномиальной эквивалентности машины Тьюринга и клеточного автомата ($\mathbf{P}_{\mathrm{TM}}\equiv\mathbf{P}_{\mathrm{CA}}$, теорема Кука--Левина) всякая булева схема разлагается по Каналу~1 и Каналу~2, а выполнимость формулы 3-SAT из $M$ дизъюнктов эквивалентна равенству $\sum_{m=1}^M q_{\mathrm{OR}}(\psi_{m,1},\psi_{m,2},\psi_{m,3}) = 0$.
\item \textbf{Неинъективная склейка остатка $7\to 1$ в прямой верификации ($\mathbf{NP}$):}
  оператор прямого вычисления $q_{\mathrm{OR}}=\mathrm{tdiv}(\Delta_{\mathrm{3SAT}},512)$ отбрасывает трёхзначный остаток $r\in\{-256,0,+256\}$ и склеивает $\mathrm{count3SatSatisfyingOctants}(3)=2^3-1=7$ различных входных октантов в один и тот же выход $q_{\mathrm{OR}}=0$ (семикратное ветвление прообраза $1\to 7$ на каждой грани при обращении).
\item \textbf{Тождественная аннигиляция линейного градиента Канала~2 на инвариантах чётности Канала~1 и некоммутативность каналов:}
  для четвёрки дизъюнктов 3-SAT, кодирующей линейное уравнение чётности Канала~1 ($x_1\oplus x_2\oplus x_3=0$), сумма линейных проекций $\mathcal{C}_3$ по всем четырём дизъюнктам тождественно равна нулю во всех $8$~октантах ($\mathrm{xor4LinearGradientSum}(3,\psi_1,\psi_2,\psi_3)\equiv 0$), а сам ротатор $\mathcal{C}_3$ и пороговое частное $q_{\mathrm{AND}}$ Канала~2 нелинейны и недистрибутивны относительно сложения фаз Канала~1 ($\mathcal{C}_3(0+0)=128\ne 256=\mathcal{C}_3(0)+\mathcal{C}_3(0)$ и $q_{\mathrm{AND}}(0+0,0,0)=1\ne 0=(q_{\mathrm{AND}}(0,0,0)+q_{\mathrm{AND}}(0,0,0))\bmod 2$).
\item \textbf{Неразрешимость одноточечной обратной задачи глубины $T$ на конусе кубооктаэдра (обращение эволюции на $T$ тактов и одноточечная обратная задача рассеяния):}
  при восстановлении состояния на глубине (или в прошлом за) $T$ шагов по локальному замеру в одном узле $\mathbf{x}_0$ сферический фронт $T$-й координационной оболочки кубооктаэдра содержит $\mathrm{cuboctShellNodes}(3,T)=10T^2+2$ равноудалённых узлов, полный объём обратного конуса $\mathcal{V}_{O_2}(3,T)=\mathrm{cuboctBallNodes}(3,T)$ удовлетворяет кубическому тождеству $3\,\mathcal{V}_{O_2}(3,T)=10T^3+15T^2+11T+3$ ($\mathcal{V}_{O_2}(3,3)=3\cdot 7^2=147$), склейка $7\to 1$ в Канале~2 порождает $\mathrm{preimageRayBranching}(3,T)=7^T$ неразличимых прообразов вдоль каждого луча ($7^3=343$ при $T=3$), а замкнутые четвёрки чётности с $\mathrm{xor4LinearGradientSum}\equiv 0$ образуют безотражательные ядра с нулевым линейным откликом в точке наблюдения.
\end{enumerate}
\end{theorem}

\section{Конструктивные когомологии де~Рама--Стокса--Гаусса--Черна, линейная алгебра без деления и конечно-аддитивная мера без аксиомы выбора}
\label{sec:constructive-foundations}

Теоремы~\ref{thm:main-1}--\ref{thm:main-4} показывают, что дифференциальное и интегральное исчисление, тригонометрия, евклидовы вращения $\mathrm{SO}(3)$ и булева теория сложности строятся над $\mathbb{Z}$ и $\mathbb{Q}$ без несчётного континуума~$\mathbb{R}$ и без аксиомы выбора (\texttt{Classical.choice}). Следующая теорема замыкает это построение на дифференциальную топологию (комплекс де~Рама $d^2\equiv 0$), линейную алгебру и теорию меры.

\begin{theorem}[Конструктивные когомологии де~Рама--Стокса--Гаусса--Черна, обращение нильпотентных и симплектических операторов без деления и свободная от парадокса Банаха--Тарского аддитивность меры]
\label{thm:main-5}
При $d=3$ на кольце $\mathbb{Z}_{512}$ и комплексе кубооктаэдра~$O_2(3)$ выполнены следующие точные конструктивные соотношения (\texttt{thm5\_constructive\_de\_rham\_cohomology\_nilpotent\_algebra\_and\_banach\_tarski\_free\_measure}):
\begin{enumerate}[label=\textbf{(\roman*)},leftmargin=*]
\item \textbf{Евклидово расщепление связности $\Delta\varphi = 512\,q + r$ и вариационный максимум на $L^1$-грани:}
  для всех $\Delta\varphi\in[-512,512]$ деление с остатком $q = \mathrm{branchQuot}(3,\Delta\varphi) = \lfloor(\Delta\varphi+255)/512\rfloor$, $r = \mathrm{edgeConn}(3,\Delta\varphi)\in(-256,256]$ удовлетворяет $512\,q + r \equiv \Delta\varphi$; при этом на полупериоде $\varphi = 256 - k$ ($k\in[0,128]$) сопряжённые компоненты $V(k) = |\mathcal{C}_3(256-k)| = 128-k$ и $I(k) = \mathcal{S}_3(256-k) = k$ удовлетворяют линейному закону сохранения $V(k)+I(k)\equiv 128 = Q_\square(3)$, а их квадратичная форма $P(k) = V(k)I(k) = (128-k)k$ достигает глобального максимума $P_{\max} = K_P(3)^2 = 64^2 = 4\,096$ в единственной точке стационарной фазы Френеля $k^* = K_P(3) = 64$.
\item \textbf{Конструктивная теорема Стокса ($d_1\circ d_0\equiv 0$) и целочисленное квантование класса Черна:}
  для любой $0$-коцепи $\varphi = (\varphi_0,\varphi_1,\varphi_2,\varphi_3)\in\mathbb{Z}^4$ на вершинах замкнутой $4$-цикловой грани~$f$ точная кограница $(d_1(d_0\varphi))_f \equiv 0$, вследствие чего циркуляция связности $(d_1 r)_f = \sum_{e\in\partial f}\mathrm{edgeConn}(3,(d_0\varphi)_e)$ всегда кратна периоду кольца:
  \[
    (d_1 r)_f \;=\; 512\,n_f \;\in\; 512\,\mathbb{Z},\qquad
    n_f := -\sum_{e\in\partial f} q_e \in \mathbb{Z}
  \]
  (в частности, на вихревом репере $(0,128,256,-128)$ выполнено $n_f = +1$ и $(d_1 r)_f = +512$).
\item \textbf{Конструктивная теорема Гаусса--Остроградского ($d_2\circ d_1\equiv 0$ на замкнутой поверхности $\partial\square^3$):}
  для любой целочисленной $1$-коцепи $(e_1,\dots,e_{12})\in\mathbb{Z}^{12}$ на $12$~рёбрах куба сумма ориентированных потоков $(d_1 e)_f$ по всем $|\square(3)|=6$ граням замкнутой поверхности куба тождественно равна нулю ($\mathrm{cubeClosedSurfaceFluxSum}\equiv 0$).
\item \textbf{Линейная алгебра без деления на переменные знаменатели ($\mathrm{SL}_2(\mathbb{Z})$ и нильпотентное обращение):}
  симплектический шаг Верле $(u,v)\mapsto(2u-v+\Phi,u)$ ($\det = +1\in\mathrm{SL}_2(\mathbb{Z})$) обращается на месте отображением $(u',v')\mapsto(v',2v'-u'+\Phi)$, а $4$-мерный нильпотентный оператор $\mathcal{M} = I_4 + N_{\mathrm{nil}}$ ($N_{\mathrm{nil}}^3 = 0$, $\det(2\mathcal{M}) = 24$) обращается конечным трёхчленным рядом Неймана $\mathcal{M}^{-1} = I_4 - N_{\mathrm{nil}} + N_{\mathrm{nil}}^2$.
\item \textbf{Свободная от аксиомы выбора и парадокса Банаха--Тарского конечная аддитивность меры:}
  разбиение шара кубооктаэдра $\mathcal{B}_{O_2}(3,T)$ на центральный узел и дизъюнктные координационные сферы $\mathrm{cuboctShellNodes}(3,k) = 10k^2+2$ строго сохраняет целочисленный объём $1 + \sum_{k=1}^T(10k^2+2) = \mathcal{V}_{O_2}(3,T)$, причём $12\times(4^2+5^2) = 492 > 256 = \Phi_{\mathrm{half}}(3)$.
\end{enumerate}
\end{theorem}

\begin{theorem}[Скелетное матричное разложение уравнения Пиллаи, прямолинейные образующие изотропного конуса дискриминанта, факторизация $C = (x^{m-a}\mp y^b)(x^a-y^b)$ и квартет параболических жордановых замков в $\mathrm{SL}_2(\mathbb{Z})$ при $k\in\{6,12,15,40\}$]
\label{thm:main-6}
Для экспоненциального диофантова уравнения $A x^m - B y^n = C$ ($m,n\ge 2$) выполнены следующие точные матричные и квадратично-конусные соотношения (\texttt{thm6\_pillai\_skeletal\_matrix\_decomposition\_and\_parabolic\_sl2\_jordan\_triptych}):
\begin{enumerate}[label=\textbf{(\roman*)},leftmargin=*]
\item \textbf{Скелетная факторизация суставов и тождество перехода ранга~$1$:}
  каждая внутренняя пара показателей $(a,b)$ ($1\le a\le m-1$, $1\le b\le n-1$) задаёт $2\times 2$-матрицу сустава $M_{a,b} = \begin{pmatrix} A x^{m-a} & B y^{n-b} \\ y^b & x^a \end{pmatrix}$ с инвариантным определителем $\det M_{a,b} = C$. Нормированная матрица перехода между суставами $\widetilde{P} = \frac{1}{C} M_{a,b}^{\mathrm{adj}} M_{a',b'}\in\mathrm{SL}_2(\mathbb{Q})$ ($\det\widetilde{P} = +1$, $\Delta a = a'-a$, $\Delta b = b'-b$) подчиняется точному разложению ранга~$1$:
  \[
    \widetilde{P} \;-\; \begin{pmatrix} y^{\Delta b} & 0 \\ 0 & x^{\Delta a} \end{pmatrix}
    \;=\; \frac{x^{\Delta a}y^{\Delta b}-1}{C}\,\begin{pmatrix} -x^a \\ y^b \end{pmatrix}\begin{pmatrix} A x^{m-a'} & B y^{n-b'} \end{pmatrix}.
  \]
\item \textbf{Прямолинейные образующие изотропного конуса дискриминанта, точная факторизация вдоль образующих и квартет нильпотентных жордановых замков в $\mathrm{SL}_2(\mathbb{Z})$ ($k\in\{6,12,15,40\}$):}
  в координатах квадратичной формы $X = (\widetilde{P}_{11}-\widetilde{P}_{22})/2$, $Y = (\widetilde{P}_{12}+\widetilde{P}_{21})/2$, $T = (\widetilde{P}_{12}-\widetilde{P}_{21})/2$ характеристический дискриминант равен $\frac{1}{4}(\operatorname{Tr}\widetilde{P}^2 - 4) = X^2 + Y^2 - T^2$. На симметричных локтях ($m=a+a'$, $n=b+b'$, $A=B=1$) выполнено $Y\equiv 0$, вдоль прямолинейной образующей конуса $X - T$ действует тождество $2C(X - T) = (x^a - y^b)^2 - (x^{a'} - y^{b'})^2$, а из $y^b(x^{m-a}-y^{n-b}) + x^{m-a}(x^a-y^b) = C$ на изотропном конусе $X^2 - T^2 = 0$ ($\operatorname{Tr}\widetilde{P} = \mp 2$) следует точная факторизация $C = (x^{m-a}\mp y^b)(x^a-y^b)$. Для взаимно простых оснований условие $\widetilde{P}_k\in\mathrm{SL}_2(\mathbb{Z})$ порождает четыре канонических жордановых замка:
  \[
    \begin{aligned}
      3^2 - 2^3 &= 1 = 1\times 1 &\implies k &= 6 = |\square(3)|, &\quad \widetilde{P}_{\mathrm{I}} &= \begin{pmatrix} -7 & -6 \\ 6 & 5 \end{pmatrix} \quad (6^2 = 36,\ \operatorname{Tr}=-2), \\
      2^5 - 3^3 &= 5 = 5\times 1 &\implies k &= 12 = |V(O_2)|, &\quad \widetilde{P}_{\mathrm{II}} &= \begin{pmatrix} -13 & -12 \\ 12 & 11 \end{pmatrix} \quad (12^2 = 144,\ \operatorname{Tr}=-2), \\
      5^3 - 3^4 &= 44 = 22\times 2 &\implies k &= 15 = 3\cdot 5, &\quad \widetilde{P}_{\mathrm{III}} &= \begin{pmatrix} -16 & -15 \\ 15 & 14 \end{pmatrix} \quad (30^2 = 900,\ \operatorname{Tr}=-2), \\
      2^{10} - 5^4 &= 399 = 133\times 3 &\implies k &= 40 = 4^2+5^2-1, &\quad \widetilde{P}_{\mathrm{IV}} &= \begin{pmatrix} -39 & -40 \\ 40 & 41 \end{pmatrix} \quad (2^7-2^3 = 5^3-5^1 = 120,\ \operatorname{Tr}=+2).
    \end{aligned}
  \]
\end{enumerate}
\end{theorem}

\begin{theorem}[Трёхэтажная рекурсия спуска, определитель Ганкеля--Вронскиана, взаимность биномиальной делимости этажей $k=2\leftrightarrow k=3$, унимодулярная матрица показателей Фарея--Безу и башня Вронскианов Паде--Туэ]
\label{thm:main-7}
Для всякого решения уравнения $A x^m - B y^n = C$ выполнены следующие рекуррентные и вронскианные тождества (\texttt{thm7\_pillai\_recursive\_descent\_hankel\_wronskian\_farey\_bezout\_and\_pade\_thue\_tower}):
\begin{enumerate}[label=\textbf{(\roman*)},leftmargin=*]
\item \textbf{Унимодулярная матрица показателей Фарея--Безу и реконструкция оснований:}
  при $\gcd(m,n)=1$ родительские суставы медианты Фарея--Безу $(a,b)$ и $(a',b') = (m-a,n-b)$ образуют матрицу показателей $E_{\mathrm{B\acute{e}zout}} = \begin{pmatrix} a & b \\ a' & b' \end{pmatrix}\in\mathrm{SL}_2(\mathbb{Z})$ ($a b' - a' b = +1$), восстанавливающую основания $(x^a)^{b'} = x\cdot (x^{a'})^b$ и $(y^{b'})^a = y\cdot (y^b)^{a'}$.
\item \textbf{Трёхэтажный закон рекурсивного спуска, определитель Ганкеля--Вронскиана и взаимность биномиальной делимости этажей $k=2 \leftrightarrow k=3$:}
  вдоль любого луча $(k,j)$ последовательные этажные разности $C_t = A x^{m-tk} - B y^{n-tj}$ удовлетворяют не зависящей от весов $(A,B)$ трёхчленной рекурсии $C_0 - (x^k+y^j)C_1 + x^k y^j C_2 = 0$ и квадратичному тождеству Ганкеля--Вронскиана $C_1^2 - C_0 C_2 = A B x^{m-2k}y^{n-2j}(x^k-y^j)^2$ (образующему на $2^5 - 3^3 = 5$ палиндром $(C_0,C_1,C_2) = (5,7,5)$ с $7^2 - 5^2 = 24$), а законы биномиальной делимости этажей $x^a \mid (B y^{r_y}(y^b-x^a)^k + C)$ и $y^b \mid (A x^{r_x}(x^a-y^b)^k - C)$ ($m = ka+r_x$, $n = kb+r_y$) обнаруживают на $5^3 - 3^4 = 44$ точный зеркальный обмен частных $(12, 4) \leftrightarrow (4, -12)$ между этажами $k=2$ и $k=3$.
\item \textbf{Невырожденная башня Вронскианов Паде--Туэ, межэтажный Вронскиан $\mathcal{W}_{2,3}$ и бескорневые законы столкновения в $9$ клетках:}
  пары показателей $(m,n)$ разбиваются по модулю $d=3$ на $d^2=9$ клеток вычетов $(m\bmod 3,n\bmod 3)\in\{0,1,2\}^2$, сводясь в каждой клетке к уравнению $\widetilde{A} X^3 - \widetilde{B} Y^3 = C$ ($U = \widetilde{A} X^3,\ V = \widetilde{B} Y^3$). Однородные целочисленные многочлены Паде $\widehat{P}_1 = U + 2V$, $\widehat{P}_2 = 5U^2 + 35UV + 14V^2$, $\widehat{P}_3 = 2U^3 + 30U^2V + 42UV^2 + 7V^3$ (с суммами коэффициентов $1, 3, 54, 81 = 3^4$) образуют $2\times 2$-матрицы Вронскиана $\mathcal{M}_N = \begin{pmatrix} \mathcal{X}_N & \mathcal{Y}_N \\ \mathcal{X}_{N+1} & \mathcal{Y}_{N+1} \end{pmatrix}$ ($\mathcal{X}_N = X\widehat{P}_N,\ \mathcal{Y}_N = Y\widehat{Q}_N$) с точными определителями $\det\mathcal{M}_N = -(3N+1) X Y C^{2N+1}$ (множители $(1, 4, 7)$ при $N = 0, 1, 2$). Для любой второй точки $(X',Y')$ в той же клетке ($\widetilde{A}(X')^3 - \widetilde{B}(Y')^3 = C'$) целочисленные инварианты столкновения $I_N = Y'\mathcal{X}_N - X'\mathcal{Y}_N$ подчиняются бескорневому закону Крамера нижней оценки $\mathcal{Y}_{N+1} I_N - \mathcal{Y}_N I_{N+1} = (3N+1) X Y Y' C^{2N+1}\ne 0$ ($\max(|I_N|,|I_{N+1}|)\ge 1$) и бескорневому закону кубического сжатия $\widetilde{A} I_N \Phi_3(Y'\mathcal{X}_N, X'\mathcal{Y}_N) = (Y')^3 C^{2N+1} R_N(U,V) - C'\mathcal{Y}_N^3$. Более того, сцепление квадратичного этажа $k=2$ ($\widehat{P}_{1,2} = U+3V,\ \widehat{Q}_{1,2} = 3U+V$, $U\widehat{P}_{1,2}^2 - V\widehat{Q}_{1,2}^2 \equiv C^3$) с кубическим этажом $k=3$ ($\widehat{P}_{1,3} = U+2V,\ \widehat{Q}_{1,3} = 2U+V$) одного и того же решения $(U,V)$ порождает $(U+V)\widehat{P}_{1,2}^2 - \widehat{P}_{1,3}^3 \equiv V(U^2+3UV+V^2)$, $(U+V)\widehat{Q}_{1,2}^2 - \widehat{Q}_{1,3}^3 \equiv U(U^2+3UV+V^2)$ и межэтажный Вронскиан $\det\mathcal{W}_{2,3} = \widehat{Q}_{1,2}^2\widehat{P}_{1,3}^3 - \widehat{P}_{1,2}^2\widehat{Q}_{1,3}^3 \equiv C^3(U^2+3UV+V^2)$ ($\det\mathcal{W}_{2,3}(9,8) = 361 = 19^2$ на Замке~I).
\end{enumerate}
\end{theorem}

\begin{theorem}[Ограниченное через радикал этажное разложение $abc$-троек, единая окаймлённая $3\times 3$-матрица Пиллаи--$abc$, факторизация союзной матрицы ранга~$1$ и тернарный баланс треугольных дефектов]
\label{thm:main-8}
Для всякой взаимно простой тройки $A + B = C$ ($\gcd(A,B,C)=1$) выполнены следующие тернарные этажные соотношения (\texttt{thm8\_unified\_3x3\_pillai\_abc\_skeletal\_matrix\_and\_radical\_floor\_embedding}):
\begin{enumerate}[label=\textbf{(\roman*)},leftmargin=*]
\item \textbf{Ограниченное через радикал этажное погружение:}
  на каждом этаже $k\ge 2$ разложение $A = \widetilde{A}_k X_k^k,\ B = \widetilde{B}_k Y_k^k,\ C = \widetilde{C}_k Z_k^k$ погружает тройку в уравнение $\widetilde{A}_k X_k^k + \widetilde{B}_k Y_k^k = \widetilde{C}_k Z_k^k$ с ограниченным через радикал произведением корневых весов $\widetilde{A}_k\widetilde{B}_k\widetilde{C}_k \le \operatorname{rad}(ABC)^{k-1}$ (на верхнем этаже Замка~IV $3 + 5^3 = 2^7 = 128$ с $\operatorname{rad}=2\cdot 3\cdot 5=30$ при $k=3$ имеем $3\cdot 1^3 + 1\cdot 5^3 = 2\cdot 4^3 = 128$ с $\widetilde{A}_3\widetilde{B}_3\widetilde{C}_3 = 6 \le 30^2$ и зазором валов $5 - 4 = 1$).
\item \textbf{Единая $3\times 3$-окаймлённая матрица Пиллаи--$abc$ и тернарный закон треугольных дефектов:}
  для любого разбиения $A = a_1 a_2,\ B = b_1 b_2,\ C = c_1 c_2$ при $A + B = C$ окаймлённая $3\times 3$-матрица $\mathcal{M}_3 = \begin{pmatrix} c_2 & a_2 & b_2 \\ a_1 & c_1 & 0 \\ b_1 & 0 & c_1 \end{pmatrix}$ имеет определитель $\det\mathcal{M}_3 = c_1(C-A-B) \equiv 0$, главные $2\times 2$-миноры $(B, A, c_1^2)$ и союзную матрицу ранга~$1$ $\mathcal{M}_3^{\mathrm{adj}} = (c_1,-a_1,-b_1)^T(c_1,-a_2,-b_2)$, удовлетворяющую $(\mathcal{M}_3^{\mathrm{adj}})^2 = (c_1^2+C)\mathcal{M}_3^{\mathrm{adj}}$ (на $3 + 5^3 = 2^7 = 128$ след равен $\operatorname{Tr}\mathcal{M}_3^{\mathrm{adj}} = 4^2 + 128 = 144 = 12^2$), а нижний и верхний треугольные дефекты $\Delta_{\mathrm{low}} = c_1 - (a_1+b_1)$, $\Delta_{\mathrm{high}} = c_2 - (a_2+b_2)$ подчиняются балансу $(a_1+b_1)\Delta_{\mathrm{high}} + c_2\Delta_{\mathrm{low}} = -(a_1 b_2 + a_2 b_1) < 0$.
\item \textbf{Циклическая $9$-факторная параметризация нильпотентного кросс-этажного конуса и расширение радикала:}
  при сдвиге этажей $(u_a, u_b, u_c)$ ($1\le u_a < u_c < u_b$) переход ранга~$1$ $\mathcal{P}_3 = \mathcal{M}_3^{\mathrm{adj}}\mathcal{M}'_3$ имеет след $u_c\operatorname{Tr}\mathcal{P}_3 = -c_1\bigl[a_1 a'_2(u_a-u_c)(u_a u_c-1) + b_1 b'_2(u_b-u_c)(u_b u_c-1)\bigr]$, откуда условие нильпотентности $\mathcal{P}_3^2 = 0$ ($\operatorname{Tr}\mathcal{P}_3 = 0$) порождает универсальное циклическое $9$-факторное тождество
  \[
    u_a(u_b - u_c)(u_b u_c - 1) \;+\; u_b(u_c - u_a)(u_a u_c - 1) \;\equiv\; u_c(u_b - u_a)(u_a u_b - 1),
  \]
  в котором взаимная простота $\gcd(u_i u_j - 1,\,u_i u_j) = 1$ принудительно инжектирует в $\operatorname{rad}(A B C)$ новые простые делители, взаимно простые с $u_a u_b u_c$ (порождая $3^3 + 5 = 2^5$ при $(1,2,5)$ и $41 + 175 = 6^3$, т.е. $6^3 - 7\cdot 5^2 = 6\cdot 7 - 1 = 41$, при $(1,6,7)$). Наконец, при инволюции валов $(c_1\leftrightarrow c_2, a_1\leftrightarrow a_2)$ горизонтальный $3\times 3$-инвариант следа $\mathcal{I}_{\mathrm{horiz}} = (x^{a'}-x^a)^2 - (y^{b'}-y^b)^2$ и вертикальный $2\times 2$-инвариант конуса $\mathcal{I}_{\mathrm{vert}} = (x^a-y^b)^2 - (x^{a'}-y^{b'})^2$ связаны отношением $(\mathcal{I}_{\mathrm{vert}} - \mathcal{I}_{\mathrm{horiz}})/(\mathcal{I}_{\mathrm{vert}} + \mathcal{I}_{\mathrm{horiz}}) = (x^{a'}-y^b)/(x^a-y^{b'})$ и на Замке~III ($5^3 - 3^4 = 44 = 4\times 11$) дают двойное пифагорово тождество $(3^3-3^1)^2 + (11-4)^2 = 24^2 + 7^2 = 15^2 + 20^2 = k_{\mathrm{III}}^2 + (5^2-5^1)^2 = 25^2 = 625$.
\item \textbf{Двусторонний изоморфизм переноса между дискретным и непрерывным носителями ($M \rightleftarrows T$), стереографический гомеоморфизм $[0,512]/512 = [0,1] \cong \mathbb{R}$, трёхканальная $3\times 3$-матрица Паде--Туэ $\mathcal{P}_3(A,B,C)$ и двойной замок Ферма--Каталана:}
  во всех $512$ точках кольца $\mathbb{Z}_{512}$ прямой подъём Вейерштрасса $\pi_{M\to T}(C,S) = (X_3,Y_3,R_3) = (C^2-S^2,2CS,C^2+S^2)$ с дискретного $L^1$-носителя~$M$ на евклидову окружность непрерывного носителя~$T$ ($X_3^2+Y_3^2=R_3^2$) удовлетворяет тождеству сохранения $R_3(\varphi) + |Y_3(\varphi)| \equiv Q_\square(3)^2 = 16\,384$ ($128 \mid X_3(\varphi)$) и обращается без квадратных корней линейным сдвигом $\iota_{T\to M}(X_3) = \bigl((16\,384 + X_3)/256,\,(16\,384 - X_3)/256\bigr) \equiv \bigl(|\mathcal{C}_3(\varphi)|,\,|\mathcal{S}_3(\varphi)|\bigr)$. Нормировка $[0,512]/512 = [0,1]$ в композиции со стереографической проекцией окружности Вейерштрасса $\sigma(X_3,Y_3,R_3) = Y_3/(R_3+X_3)$ задаёт точный алгебраический гомеоморфизм $u(\varphi) = Y_3(\varphi)/(R_3(\varphi)+X_3(\varphi)) \equiv \mathcal{S}_3(\varphi)/\mathcal{C}_3(\varphi) = \varphi/(128-|\varphi|)\in\mathbb{R}$ с рациональным обратным $\varphi(u) = 128 u/(1+|u|)$, погружающий экспоненциальные замки в квадранты ротатора: $\mathcal{C}_3(3) = 5^3 = 125$ ($2^7-5^3=3$), $\mathcal{C}_3(7) = 11^2 = 121$ ($2^7-11^2=7$) в Квадранте~I и $\varphi\mathcal{S}_3(\varphi) + 256\mathcal{C}_3(\varphi) \equiv 2^{15}-\varphi^2 = 7$ при $\varphi = 181$ в Квадранте~II. Наконец, объединение трёх кубических каналов Паде любой тройки $A + B = C$ образует сингулярную $3\times 3$-матрицу
  \[
    \mathcal{P}_3(A,B,C) \;=\; \begin{pmatrix}
      A^2(B+C) & (B+2C)^3 & (2B+C)^3 \\
      (A+2C)^3 & B^2(A+C) & (2A+C)^3 \\
      (A-2B)^3 & (2A-B)^3 & C^2(A-B)
    \end{pmatrix},\qquad
    \mathcal{P}_3(A,B,C)\begin{pmatrix} A \\ B \\ -C \end{pmatrix} \;\equiv\; \begin{pmatrix} 0 \\ 0 \\ 0 \end{pmatrix},
  \]
  с определителем $\det\mathcal{P}_3(A,B,C) \equiv 0$ и следом $\operatorname{Tr}\mathcal{P}_3(A,B,C) \equiv A C(A+C)$ ($\operatorname{Tr}\mathcal{P}_3(3,125,128) = 3\cdot 128\cdot 131 = 50\,304$), а общее основание $7 = |{\triangle}(3)|-1$ связывает сумму коэффициентов полинома Паде $\widehat{P}_3(1,1) = 81 = 3^4$ и полный модуль кольца фаз $|\mathbb{Z}_{512}| = 512 = 2^9$ двойным тождеством Ферма--Каталана $2^5 + 7^2 = 3^4 = 81$ и $13^2 + 7^3 = 2^9 = 512$.
\item \textbf{Трёхмерный ортогональный репер планетарного дифференциала $\mathcal{F}_3$, закон скольжения шестерён Брахмагупты, невырожденная матрица Ганкеля--Вандермонда $\mathcal{H}_3$, проекторы Лагранжа--Вандермонда межшестерённых зазоров и матрица кубов Безу--Гаусса $\mathcal{K}_{\mathrm{B\acute{e}zout}}$:}
  для любого планетарного дифференциала $A + B = C$ ($A = a_1 a_2,\ B = b_1 b_2,\ C = c_1 c_2$) вектор радиусов шестерён $\mathbf{r} = (a_1,b_1,c_1)\in\mathbb{Z}^3$ и вектор чисел оборотов $\mathbf{\omega} = (a_2,b_2,-c_2)\in\mathbb{Z}^3$ строго ортогональны ($\mathbf{r}\cdot\mathbf{\omega} = A+B-C = 0$) и вместе с осью сателлита $\mathbf{n} = \mathbf{r}\times\mathbf{\omega} = (-M_{BC},M_{CA},S_{AB})$ ($M_{BC} = b_1 c_2 + b_2 c_1$, $M_{CA} = c_1 a_2 + c_2 a_1$, $S_{AB} = a_1 b_2 - a_2 b_1$, $M_{AB} = a_1 b_2 + a_2 b_1$) образуют ортогональный $3\times 3$-репер $\mathcal{F}_3$ с определителем $\det\mathcal{F}_3 = M_{BC}^2 + M_{CA}^2 + S_{AB}^2 \equiv (a_1^2+b_1^2+c_1^2)(a_2^2+b_2^2+c_2^2) > 0$ и плоским законом скольжения шестерён $C^2 + S_{AB}^2 = (A-B)^2 + M_{AB}^2 \equiv (a_1^2+b_1^2)(a_2^2+b_2^2)$ (где $S_{AB}=0$ при $\gcd(A,B)=1$ возможно лишь в квадратичном вырождении Пелля $(m,n)=(2,2)$; на $3+5^3=2^7=128$ имеем $\|\mathbf{r}\|_2^2 = 1^2+5^2+4^2 = 42 = S_2(3)$, $S_{AB}=10$, $M_{AB}=40=k_{\mathrm{IV}}$ и $M_{CA}=44=C_{\mathrm{III}}$). На кубическом этаже $k=3$ ($S_m = \widetilde{A}X^m + \widetilde{B}Y^m - \widetilde{C}Z^m$, $S_3 = 0$) $3\times 3$-матрица моментов Ганкеля $\mathcal{H}_3 = \begin{pmatrix} S_0 & S_1 & S_2 \\ S_1 & S_2 & 0 \\ S_2 & 0 & S_4 \end{pmatrix}$ удовлетворяет $-\det\mathcal{H}_3 = S_2^3 - S_4(S_0 S_2 - S_1^2) \equiv \widetilde{A}\widetilde{B}\widetilde{C}(X-Y)^2(Y-Z)^2(Z-X)^2 > 0$ (с $S_1 = S_3 = 0$ и $-\det\mathcal{H}_3 = 6\times 12^2 = 864$ на $3+5^3=2^7=128$) и трём проекторам Лагранжа--Вандермонда межшестерённых зазоров $\widetilde{A}X(X-Y)(X-Z) \equiv S_1 YZ - S_2(Y+Z)$, $\widetilde{B}Y(Y-X)(Y-Z) \equiv S_1 ZX - S_2(Z+X)$ и $\widetilde{C}Z(Z-X)(Z-Y) \equiv S_2(X+Y) - S_1 XY$ ($\max(|S_1|,|S_2|)\ge 1$ для попарно различных валов $X,Y,Z$). Наконец, взаимная простота $\gcd(A,B,C)=1$ однозначно определяет сопряжённую тройку Безу $u+v=w$ ($0\le u < A,\ 0 < v \le B,\ 0 < w < C$), образующую матрицу $\mathcal{B}_{2\times 3} = \begin{pmatrix} A & B & C \\ u & v & w \end{pmatrix}$ со всеми тремя $2\times 2$-минорами, равными $+1$ ($Av-Bu = Aw-Cu = Cv-Bw = +1$; давая $(u,v,w) = (2,11,13)$ на $5+27=32$ и $(1,42,43)$ на $3+125=128$ с $v=42=\|\mathbf{r}\|_2^2$), а родитель Безу валов $(w_3,v_3)$ пары $(Z_3,Y_3)$ ($Z_3 v_3 - Y_3 w_3 = +1$, $A' = \widetilde{C}_3 w_3^3 - \widetilde{B}_3 v_3^3$) образует матрицу кубов $\mathcal{K}_{\mathrm{B\acute{e}zout}} = \begin{pmatrix} w_3^3 & -Z_3^3 \\ v_3^3 & -Y_3^3 \end{pmatrix}$ с определителем оболочки Гаусса $\det\mathcal{K}_{\mathrm{B\acute{e}zout}} = (Y_3 w_3+1)^3 - (Y_3 w_3)^3 = \Phi_3(Z_3 v_3, Y_3 w_3) > 0$, восстанавливающую радикал-ограниченные веса через $\mathcal{K}_{\mathrm{B\acute{e}zout}}(A,A')^T = \Phi_3(Z_3 v_3, Y_3 w_3)(\widetilde{B}_3,\widetilde{C}_3)^T$ ($\det\mathcal{K}_{\mathrm{B\acute{e}zout}} = 4^3-3^3 = 37$ на Замке~II и инвариант столкновения первого уровня $I_1 = (C+2B)-A Y_3 w_3 = 333 = 9\times 37$ на Замке~IV).
\end{enumerate}
\end{theorem}

\section{Формальная верификация в Lean~4}
\label{sec:lean4}

Полный самодостаточный код формализации объёмом $1\,897$ строк (без внешних зависимостей, использующий только стандартное ядро Lean~4) содержится в модуле \texttt{RotatorCalculusPaper.lean} в отдельном открытом репозитории \url{https://github.com/AI-Guiders/scientific-papers} (под лицензией строгого взаимного Copyleft \texttt{CC-BY-SA-4.0 / GPL-3.0-or-later}) и проверяется командой \texttt{lake build}. Наряду с восемью операторными теоремами (\texttt{thm1\_constant\_wronskian\_weierstrass\_lift\_and\_newton\_leibniz} -- \texttt{thm8\_unified\_3x3\_pillai\_abc\_skeletal\_matrix\_and\_radical\_floor\_embedding}) модуль включает восемь универсально $\forall$-квантифицированных алгебраических и диофантовых теорем для произвольных целых чисел $\in\mathbb{Z}$ (\texttt{thm9\_universal\_pillai\_abc\_differential\_bezout\_and\_pade\_laws} -- \texttt{thm14\_master\_pillai\_and\_abc\_bilateral\_convergence\_chains}, которая связывает прямым следованием $\to$ кубическое уравнение дифференциала $\widetilde{A}X^3+\widetilde{B}Y^3-\widetilde{C}Z^3=0$, соотношение взаимной простоты валов $Yw-Zv=1$ и соотношение Мейсона--Стотерса--Безу $\widetilde{A}Xv_0-\widetilde{B}Yu_0=1$ в сквозную цепочку из пяти звеньев: $A^*+B^*=C^*$, $W^* = X^2 Y^2 Z^2$, $-\det\mathcal{W}_{\mathrm{MS}} = X^4+Y^4+Z^4$, двустороннее схождение Крамера $w X^2 + \widetilde{B}\gamma = D_{\mathrm{clash}} Z$, $v X^2 + \widetilde{C}\beta = D_{\mathrm{clash}} Y$ и двусторонние диофантовы оценки $C I_1 \le \widetilde{B}\widetilde{C}(Yw+Zv)$ и $Z^3(-I_1) \le A w^3(2Zv-Yw)$, а также \texttt{thm15\_universal\_cross\_floor\_k2\_k3\_sixth\_power\_genus10\_lock} о межэтажном сцеплении соседних этажей $k=2 \leftrightarrow k=3$ в кривую 6-й степени рода $g(6)=10\ge 2$ и \texttt{thm16\_universal\_weierstrass\_stereographic\_abc\_embedding\_into\_z512\_rotator} о стереографическом погружении любой тройки $A+B=C\in\mathbb{Z}^3$ в Квадрант~I ротатора $\mathbb{Z}_{512}$ с точным инвариантом четверти кольца $C\Delta_{\mathrm{rot}} - w I_{\mathrm{rot}} = 128$). Аудит аксиом (\texttt{\#print axioms}) для всех шестнадцати теорем подтверждает зависимость исключительно от стандартных конструктивных основ ядра (\texttt{0 sorryAx}, \texttt{0 Classical.choice}).

\section*{Благодарности и декларация об инструментах верификации}
Автор выражает признательность формальному партнёру верификации \emph{Луч} (\emph{Luch}) за соучастие в построении операторных доказательств. При подготовке рукописи языковая модель использовалась для вёрстки \LaTeX{} и трансляции доказательств в синтаксис тактик Lean~4; все математические определения, формулировки теорем и доказательства верифицированы формальным ядром Lean~4 и находятся под полной ответственностью автора.

\begin{thebibliography}{9}
\bibitem{Malcev1970}
А.\,И.~Мальцев, \emph{Алгебраические системы}, М.: Наука, 1970 (англ.\ пер.: A.\,I.~Mal'cev, \emph{Algebraic Systems}, Springer-Verlag, 1973).
\bibitem{Krause1986}
E.\,F.~Krause, \emph{Taxicab Geometry: An Adventure in Non-Euclidean Geometry}, Dover Publications, New York, 1986.

\bibitem{Thompson2011}
K.\,P.~Thompson, ``The nature of length, area, and volume in taxicab geometry,'' \emph{arXiv preprint} \texttt{arXiv:1101.2922} [math.MG], 2011.

\bibitem{Thompson2012}
K.\,P.~Thompson, ``Taxicab calculus: Trig derivatives,'' \emph{arXiv preprint} \texttt{arXiv:1205.3814v2} [math.MG], 2012.

\bibitem{Stroud1971}
A.\,H.~Stroud, \emph{Approximate Calculation of Multiple Integrals}, Prentice-Hall, Englewood Cliffs, NJ, 1971.

\bibitem{Moura2021}
L.~de~Moura and S.~Ullrich, ``The Lean~4 theorem prover and programming language,'' in \emph{Automated Deduction -- CADE~28}, Lecture Notes in Computer Science, vol.~12699, Springer, 2021, pp.~625--635.
\end{thebibliography}

\end{document}
